% ─────────────────────────────────────────────────────────────────────────────
% WERSJA POLSKA — ROZDZIAŁ 2 — WERSJE ROBOCZE
% Wszystkie wersje mają identyczną strukturę sekcji.
% Różnią się stylem:
%   Wersja A — matematyczna: kluczowe równania i notacja formalna
%   Wersja B — ogólna: konceptualna, przystępna, bez wzorów
%   Wersja C — precyzyjna, minimalnie matematyczna: dokładna terminologia, bez wyprowadzeń
% ─────────────────────────────────────────────────────────────────────────────


% ═════════════════════════════════════════════════════════════════════════════
% WERSJA A — matematyczna
% ═════════════════════════════════════════════════════════════════════════════

\chapter{Podstawy teoretyczne}

Niniejszy rozdział prezentuje formalne podstawy algorytmów ocenianych w~pracy.
Omówiono kolejno: podproblemy nawigacji mobilnej, sensory stosowane w~eksperymentach,
reprezentacje map, metody rejestracji chmur punktów, filtracyjną i~grafową
estymację stanu, preintegrację IMU oraz podstawy wizyjnego SLAM.
Rozdział zamykają metryki ewaluacji używane w~rozdziale~4.

\section{Nawigacja w~robotyce mobilnej}

Nawigacja autonomiczna polega na celowym przemieszczaniu się robota w~środowisku.
Klasycznie dekomponuje się ją na trzy powiązane podproblemy: zliczanie drogi
(\textit{dead reckoning}), mapowanie i~lokalizację.
SLAM łączy dwa ostatnie i~stanowi główny przedmiot niniejszej pracy.

\subsection{Zliczanie drogi}

Zliczanie drogi polega na szacowaniu pozy robota $\mathbf{x}_t \in SE(3)$
poprzez całkowanie pomiarów odometrycznych lub sterowań:
\begin{equation}
    \mathbf{x}_t = \mathbf{x}_{t-1} \oplus \mathbf{u}_t,
\end{equation}
gdzie $\mathbf{u}_t$ jest przyrostem pozy, a~$\oplus$ oznacza złożenie
w~grupie $SE(3)$.
Błędy kolejnych przyrostów kumulują się, powodując narastający dryf pozycji
bez żadnego mechanizmu korekcji.
W~kompletnych systemach SLAM zliczanie drogi pełni rolę modelu predykcji.

\subsection{Mapowanie}

Mapowanie to zadanie wyznaczenia modelu przestrzennego $\mathbf{m}$ ze zbioru
pomiarów $\mathbf{z}_{1:t}$ przy założeniu znanych poz $\mathbf{x}_{1:t}$:
\begin{equation}
    p(\mathbf{m} \mid \mathbf{z}_{1:t}, \mathbf{x}_{1:t}).
\end{equation}
Bez dokładnych estymacji pozycji robota kolejne pomiary nie mogą być
spójnie wbudowane w~mapę.

\subsection{Lokalizacja}

Lokalizacja polega na wyznaczeniu rozkładu pozy robota $p(\mathbf{x}_t)$
przy danej mapie $\mathbf{m}$ i~historii pomiarów.
Globalną lokalizację (brak apriorycznej informacji o~pozycji) rozwiązuje
m.in.\ metoda Monte Carlo, reprezentująca rozkład jako zbiór cząsteczek.
Śledzenie pozy zakłada ciągłość ruchu i~aproksymuje rozkład gaussowsko.

\subsection{SLAM}

SLAM łączy mapowanie i~lokalizację, szacując jednocześnie trajektorię
$\mathbf{x}_{1:t}$ i~mapę $\mathbf{m}$ ze strumienia pomiarów:
\begin{equation}
    p(\mathbf{x}_{1:t}, \mathbf{m} \mid \mathbf{z}_{1:t}, \mathbf{u}_{1:t}).
\end{equation}
W~podejściu \textit{online SLAM} marginalizuje się przeszłe pozy,
szacując jedynie $p(\mathbf{x}_t, \mathbf{m})$.
W~\textit{full SLAM} zachowuje się pełną trajektorię, co umożliwia
globalną optymalizację po wykryciu pętli.

\subsection{Aktywny SLAM}

Aktywny SLAM steruje ruchem robota tak, by minimalizować niepewność mapy
i~trajektorii, np.\ przez maksymalizację oczekiwanego przyrostu informacji.
W~niniejszej pracy stosowany jest pasywny SLAM na nagranych danych,
dlatego aktywny SLAM pozostaje poza zakresem rozważań.

\section{Sensory}

Zestaw sensorów determinuje, jakie dane są dostępne dla algorytmu SLAM.
W~eksperymentach stosowano cztery typy sensorów, omówione poniżej.

\subsection{LiDAR}

LiDAR wyznacza odległości metodą czasu przelotu impulsu laserowego:
\begin{equation}
    d = \frac{c \cdot \Delta t}{2},
\end{equation}
gdzie $c$ jest prędkością światła, a~$\Delta t$ --- zmierzonym czasem
przelotu.
Wynikiem każdego obrotu głowicy jest chmura punktów $\mathcal{P} \subset \mathbb{R}^3$
o~rozdzielczości kątowej rzędu $0{,}1$--$0{,}4°$ i~zasięgu do kilkudziesięciu metrów.
Przy szybkim ruchu robota punkty rejestrowane na początku i~końcu obrotu
odpowiadają różnym pozycjom głowicy, powodując zniekształcenie geometryczne
(\textit{motion distortion}).

\subsection{Kamera stereo}

Głębokość punktu $Z$ wyznaczana jest ze stereoskopii na podstawie dysparycji $d$:
\begin{equation}
    Z = \frac{f \cdot B}{d},
\end{equation}
gdzie $f$ to ogniskowa, a~$B$ --- baza (odległość między obiektywami).
Niepewność głębokości rośnie kwadratowo z~odległością.
Kamery dostarczają informacji teksturowej przydatnej w~dopasowaniu cech wizyjnych,
lecz są wrażliwe na warunki oświetleniowe.

\subsection{IMU}

Centrala inercyjna dostarcza pomiarów przyspieszenia $\tilde{\mathbf{a}} \in \mathbb{R}^3$
i~prędkości kątowej $\tilde{\boldsymbol{\omega}} \in \mathbb{R}^3$ z~częstotliwością
rzędu $200$~Hz, obarczonych biasem $\mathbf{b}$ i~szumem gaussowskim $\boldsymbol{\eta}$:
\begin{equation}
    \tilde{\mathbf{a}} = \mathbf{a} + \mathbf{b}_a + \boldsymbol{\eta}_a, \qquad
    \tilde{\boldsymbol{\omega}} = \boldsymbol{\omega} + \mathbf{b}_g + \boldsymbol{\eta}_g.
\end{equation}
Bezpośrednie całkowanie pomiarów IMU prowadzi do szybkiego dryfu pozycji,
jednak ciasne sprzężenie z~LiDAR-em lub kamerą umożliwia efektywną korekcję.

\subsection{Odometria kołowa}

Enkodery kół mierzą liczbę impulsów $n_L, n_R$ na lewym i~prawym kole.
Przyrost pozycji w~modelu różnicowym:
\begin{equation}
    \Delta x = \frac{\Delta s_L + \Delta s_R}{2}\cos\theta, \quad
    \Delta y = \frac{\Delta s_L + \Delta s_R}{2}\sin\theta, \quad
    \Delta\theta = \frac{\Delta s_R - \Delta s_L}{L},
\end{equation}
gdzie $\Delta s = n \cdot r \cdot 2\pi / N$ ($r$ --- promień koła, $N$ --- liczba
impulsów na obrót, $L$ --- rozstaw kół).
Poślizg i~błędy kalibracji powodują systematyczny dryf.

\section{Reprezentacje map}

Reprezentacja mapy determinuje koszt pamięciowy, szybkość aktualizacji
i~przydatność do planowania trasy lub inspekcji.

\subsection{Siatka zajętości}

Siatka zajętości dzieli przestrzeń na komórki o~boku $r$, z~których każda
przechowuje prawdopodobieństwo zajętości w~postaci log-odds:
\begin{equation}
    l_t = l_{t-1} + \log\frac{p(occ \mid z_t)}{1 - p(occ \mid z_t)}
                 - \log\frac{p_0}{1 - p_0},
\end{equation}
gdzie $p_0$ jest apriorycznym prawdopodobieństwem.
Aktualizacja inkrementalna jest szybka; reprezentacja wspiera bezpośrednio
planowanie ścieżek metodą A* lub Dijkstry.

\subsection{Chmura punktów 3D}

Chmura punktów to nieuporządkowany zbiór wektorów $\{(x_i, y_i, z_i)\}_{i=1}^N$,
opcjonalnie rozszerzony o~intensywność lub kolor.
Mapa globalna powstaje przez akumulację kolejnych skanów przy użyciu estymowanych
poz.
Reprezentacja zachowuje pełną geometrię 3D, lecz nie koduje explicite
wolnej przestrzeni i~wymaga dużej pamięci dla długich trajektorii.

\section{Rejestracja chmur punktów}

Rejestracja wyznacza transformację sztywną $(\mathbf{R}, \mathbf{t}) \in SE(3)$,
która minimalizuje odległość między dwoma zachodzącymi na siebie chmurami punktów.

\subsection{Iterative Closest Point (ICP)}

ICP minimalizuje punkt-do-punktu odległość między chmurą źródłową $\mathcal{S}$
a~docelową $\mathcal{T}$:
\begin{equation}
    (\mathbf{R}^*, \mathbf{t}^*) = \argmin_{\mathbf{R},\mathbf{t}}
    \sum_{i} \| \mathbf{R}\mathbf{s}_i + \mathbf{t} - \mathbf{t}_i \|^2,
\end{equation}
naprzemiennie przypisując korespondencje (najbliższy sąsiad w~k-d drzewie)
i~rozwiązując optymalizację metodą SVD.
Algorytm jest wrażliwy na inicjalizację i~wymaga dostatecznego pokrycia chmur.

\subsection{Point-to-Plane ICP}

Zamiast odległości punkt-do-punktu minimalizowana jest odległość punktu źródłowego
od płaszczyzny stycznej do powierzchni docelowej w~punkcie korespondencji:
\begin{equation}
    \sum_{i} \bigl((\mathbf{R}\mathbf{s}_i + \mathbf{t} - \mathbf{t}_i)
    \cdot \hat{\mathbf{n}}_i\bigr)^2,
\end{equation}
gdzie $\hat{\mathbf{n}}_i$ jest normalną powierzchni.
Metryka ta jest bardziej informatywna na gładkich powierzchniach i~przyspiesza
zbieżność.

\subsection{Generalized ICP (GICP)}

GICP modeluje każdy punkt jako rozkład gaussowski nad lokalną geometrią powierzchni.
Minimalizowane jest odchylenie Mahalanobisa między dopasowanymi rozkładami:
\begin{equation}
    \sum_{i} \mathbf{d}_i^\top
    \bigl(\mathbf{\Sigma}_i^{\mathcal{T}} + \mathbf{R}\mathbf{\Sigma}_i^{\mathcal{S}}\mathbf{R}^\top\bigr)^{-1}
    \mathbf{d}_i,
\end{equation}
gdzie $\mathbf{d}_i = \mathbf{t}_i - \mathbf{R}\mathbf{s}_i - \mathbf{t}$.
GICP łączy punkt-do-punktu i~punkt-do-płaszczyzny w~jednym probabilistycznym
sformułowaniu i~jest domyślnym wariantem stosowanym przez kilka algorytmów
ocenianych w~tej pracy.

\section{Filtracyjna estymacja stanu}

Metody filtracyjne reprezentują rozkład prawdopodobieństwa stanu robota
i~aktualizują go rekurencyjnie w~cyklu predykcja--korekcja.

\subsection{Rozszerzony filtr Kalmana (EKF)}

EKF zakłada gaussowski rozkład stanu $p(\mathbf{x}_t) = \mathcal{N}(\boldsymbol{\mu}_t, \boldsymbol{\Sigma}_t)$
i~linearyzuje modele ruchu $f$ i~obserwacji $h$ przez rozwinięcie Taylora.
Predykcja:
\begin{equation}
    \bar{\boldsymbol{\mu}}_t = f(\boldsymbol{\mu}_{t-1}, \mathbf{u}_t), \qquad
    \bar{\boldsymbol{\Sigma}}_t = \mathbf{F}_t \boldsymbol{\Sigma}_{t-1} \mathbf{F}_t^\top + \mathbf{Q}_t,
\end{equation}
korekcja:
\begin{equation}
    \mathbf{K}_t = \bar{\boldsymbol{\Sigma}}_t \mathbf{H}_t^\top
    (\mathbf{H}_t \bar{\boldsymbol{\Sigma}}_t \mathbf{H}_t^\top + \mathbf{R}_t)^{-1}, \qquad
    \boldsymbol{\mu}_t = \bar{\boldsymbol{\mu}}_t + \mathbf{K}_t(\mathbf{z}_t - h(\bar{\boldsymbol{\mu}}_t)),
\end{equation}
gdzie $\mathbf{F}_t$, $\mathbf{H}_t$ są jakobianami, $\mathbf{Q}_t$, $\mathbf{R}_t$ ---
macierzami kowariancji szumu.

\subsection{Filtr cząsteczkowy Rao-Blackwellized}

Poza wyznacza się próbkowaniem cząsteczkowym; mapa warunkowa jest analityczna.
Wagi cząsteczek po korekcji:
\begin{equation}
    w_t^{[k]} \propto \frac{p(\mathbf{z}_t \mid \mathbf{x}_t^{[k]}, \mathbf{m}^{[k]}) \,
    p(\mathbf{x}_t^{[k]} \mid \mathbf{x}_{t-1}^{[k]}, \mathbf{u}_t)}
    {\pi(\mathbf{x}_t^{[k]} \mid \mathbf{x}_{t-1}^{[k]}, \mathbf{u}_t, \mathbf{z}_t)}.
\end{equation}
Po normalizacji i~resampingu cząsteczki o~niskich wagach są zastępowane.
Algorytm GMapping, stosowany w~niniejszej pracy, opiera się na tym sformułowaniu.

\section{Grafowy SLAM}

W~grafowym SLAM trajektoria robota i~mapa są reprezentowane jako graf,
którego węzły odpowiadają pozom, a~krawędzie --- wiązaniom z~modeli odometrii
lub dopasowania skanów.

\subsection{Grafy pozycji}

Optymalizacja grafu minimalizuje sumaryczny ważony błąd wiązań:
\begin{equation}
    \mathbf{x}^* = \argmin_{\mathbf{x}} \sum_{(i,j) \in \mathcal{E}}
    \mathbf{e}_{ij}(\mathbf{x}_i, \mathbf{x}_j)^\top
    \boldsymbol{\Omega}_{ij}
    \mathbf{e}_{ij}(\mathbf{x}_i, \mathbf{x}_j),
\end{equation}
gdzie $\mathbf{e}_{ij}$ jest błędem wiązania, a~$\boldsymbol{\Omega}_{ij}$
--- macierzą informacyjną (odwrotność kowariancji).
Nieliniowy problem rozwiązywany jest iteracyjnie metodą Gaussa-Newtona
lub Levenberga-Marquardta (biblioteki g2o, GTSAM, Ceres).

\subsection{Grafy czynnikowe}

Grafy czynnikowe uogólniają grafy pozycji: węzłami mogą być dowolne zmienne
(pozy, punkty charakterystyczne, biasy IMU), a~krawędziami --- czynniki
(funkcje gęstości prawdopodobieństwa nad połączonymi zmiennymi).
Algorytmy LIO-SAM i~MOLA stosują ten formalizm.

\subsection{Detekcja i~zamykanie pętli}

Zamykanie pętli dodaje długodystansowe krawędzie do grafu, gdy robot rozpoznaje
wcześniej odwiedzone miejsce.
Dla LiDAR-u stosuje się deskryptory Scan Context lub kandydatów w~promieniu
szukania z~weryfikacją przez ICP.
Bez zamykania pętli błąd pozycji rośnie nieograniczenie wzdłuż trajektorii.

\section{Preintegracja IMU i~deskewing}

\textit{Deskewing} koryguje zniekształcenie chmury punktów wynikające z~ruchu
robota podczas obrotu głowicy LiDAR-u.
Każdy punkt $\mathbf{p}_k$ jest transformowany do wspólnej chwili $t_0$
z~użyciem interpolowanej transformacji z~IMU:
\begin{equation}
    \mathbf{p}_k^{t_0} = \mathbf{T}(t_0, t_k) \cdot \mathbf{p}_k.
\end{equation}
Preintegracja IMU łączy wszystkie pomiary między klatkami kluczowymi $i$ i~$j$
w~jedno wiązanie $\Delta\mathbf{R}_{ij}$, $\Delta\mathbf{v}_{ij}$, $\Delta\mathbf{p}_{ij}$,
bez konieczności ponownego całkowania przy każdej iteracji optymalizatora.
Zmienne integrowane w~przestrzeni $SO(3)$ zachowują strukturę grupy Liego.

\section{Podstawy wizyjnego SLAM}

Wizyjny SLAM używa obrazów kamery jako głównego źródła informacji.
W~niniejszej pracy oceniane są systemy oparte na cechach wizyjnych (sparse).

\subsection{Stereoskopia i~triangulacja}

Głębokość punktu $\mathbf{p} = (X, Y, Z)^\top$ wyznaczana jest ze wzoru
$Z = fB/d$ (por.\ sekcja~\ref{sec:sensors}).
Niepewność głębokości $\sigma_Z \propto Z^2 / (fB)$ rośnie kwadratowo
z~odległością, ograniczając wiarygodność punktów dalekich.

\subsection{Detekcja i~opis cech}

Detektor ORB wyznacza zorientowane narożniki FAST i~oblicza binarny deskryptor
BRIEF z~rotacyjną niezmienniczością.
Dopasowanie par deskryptorów metodą brute-force z~testem Lowe'a:
\begin{equation}
    \frac{d_1}{d_2} < \theta,
\end{equation}
gdzie $d_1$, $d_2$ to odległości do pierwszego i~drugiego kandydata.

\subsection{Odometria wizyjna}

Wzajemna poza kamer wyznaczana jest z~dekompozycji macierzy esencjalnej
lub przez minimalizację błędu reprojekcji w~lokalnym bundle adjustment:
\begin{equation}
    \min_{\{\mathbf{T}_i\}, \{\mathbf{P}_j\}}
    \sum_{i,j} \rho\bigl(\| \pi(\mathbf{T}_i \mathbf{P}_j) - \mathbf{u}_{ij} \|^2\bigr),
\end{equation}
gdzie $\pi$ jest funkcją projekcji, a~$\rho$ --- kosztem Huber.

\subsection{Rozpoznawanie miejsc metodą Bag-of-Words}

Słownik wizualny $\mathcal{W}$ budowany jest metodą k-means na deskryptorach
ze zbioru treningowego.
Każdy obraz reprezentowany jest histogramem $\mathbf{v} \in \mathbb{R}^{|\mathcal{W}|}$
(z~ważeniem TF-IDF).
Kandydatem zamknięcia pętli jest obraz o~najwyższym podobieństwie kosinusowym
do bieżącego, weryfikowany geometrycznie przez RANSAC.

\section{Metryki ewaluacji}

\subsection{Wyrównanie Umeyamy}

Przed obliczeniem błędów trajektoria estymowana wyrównywana jest do referencji
metodą Umeyamy: minimalizowany jest błąd RMSE między odpowiadającymi pozami
przez optymalizację transformacji podobieństwa (SVD).
Dla systemów metrycznych skala jest zablokowana.

\subsection{Bezwzględny błąd trajektorii (ATE)}

\begin{equation}
    \mathrm{ATE} = \sqrt{\frac{1}{N}\sum_{i=1}^{N}
    \| \mathbf{t}_i^{\mathrm{est}} - \mathbf{t}_i^{\mathrm{ref}} \|^2},
\end{equation}
gdzie $\mathbf{t}_i$ jest translacyjną częścią pozy.
ATE mierzy globalną spójność trajektorii.

\subsection{Względny błąd pozycji (RPE)}

\begin{equation}
    \mathrm{RPE} = \sqrt{\frac{1}{M}\sum_{k=1}^{M}
    \| \mathbf{t}(\mathbf{Q}_k^{-1}\mathbf{P}_k) \|^2},
\end{equation}
gdzie $\mathbf{P}_k = (\mathbf{T}_i^{\mathrm{ref}})^{-1}\mathbf{T}_j^{\mathrm{ref}}$,
$\mathbf{Q}_k = (\mathbf{T}_i^{\mathrm{est}})^{-1}\mathbf{T}_j^{\mathrm{est}}$
dla par $(i,j)$ odległych o~stały krok.
RPE mierzy lokalną jakość odometryczną.


% ═════════════════════════════════════════════════════════════════════════════
% WERSJA B — ogólna, konceptualna, bez wzorów
% ═════════════════════════════════════════════════════════════════════════════

\chapter{Podstawy teoretyczne}

Niniejszy rozdział omawia pojęcia i~metody, których znajomość jest niezbędna
do zrozumienia algorytmów ocenianych w~pracy.
Przedstawiono kolejno podproblemy nawigacji mobilnej, sensory użyte
w~eksperymentach, sposoby reprezentowania map, metody wyrównywania chmur punktów,
filtracyjne i~grafowe podejścia do SLAM, preintegrację IMU oraz podstawy
wizyjnego SLAM.
Na końcu opisano metryki służące do oceny dokładności trajektorii i~map.

\section{Nawigacja w~robotyce mobilnej}

Nawigacja autonomiczna to zdolność robota do celowego przemieszczania się
w~środowisku bez ingerencji człowieka.
Tradycyjnie rozkłada się ją na trzy wzajemnie powiązane zadania:
zliczanie drogi, mapowanie i~lokalizację.
SLAM łączy mapowanie i~lokalizację w~jednym, spójnym procesie i~stanowi
fundament wszystkich algorytmów omawianych w~tej pracy.

\subsection{Zliczanie drogi}

Zliczanie drogi to najprostsza forma estymacji pozycji: robot śledzi
własny ruch, sumując kolejne przyrosty przemieszczenia odczytane z~czujników
ruchu.
Metoda nie wymaga żadnej znajomości otoczenia, lecz jej słabością jest
narastający błąd --- nawet małe przekłamanie każdego pomiaru kumuluje się
z~czasem, prowadząc do coraz większego odchylenia od rzeczywistej pozycji.
W~bardziej rozbudowanych systemach zliczanie drogi służy jako szybki
i~tani punkt startowy dla dokładniejszych metod lokalizacji.

\subsection{Mapowanie}

Mapowanie polega na budowaniu przestrzennego modelu otoczenia na podstawie
odczytów sensorycznych.
Kluczowym założeniem jest znajomość pozycji robota w~chwili rejestrowania
każdego pomiaru --- bez niej kolejne obserwacje nie mogą być spójnie złożone
w~całość.
Wynikiem mapowania może być dwuwymiarowa siatka zajętości lub trójwymiarowa
chmura punktów, w~zależności od stosowanych sensorów i~algorytmu.

\subsection{Lokalizacja}

Lokalizacja to odwrotność mapowania: mając gotową mapę, robot wyznacza
własną pozycję względem niej na podstawie bieżących pomiarów sensorycznych.
Wyróżnia się lokalizację globalną, w~której robot nie posiada żadnej
wstępnej wiedzy o~swojej pozycji, oraz śledzenie pozy, które zakłada
ciągłość ruchu i~korzysta z~poprzedniej estymacji jako punktu startowego.

\subsection{SLAM}

SLAM rozwiązuje jednocześnie problem mapowania i~lokalizacji, nie zakładając
ani znajomości mapy, ani dokładnej pozycji robota.
Algorytm buduje mapę otoczenia i~estymuje trajektorię robota w~jednym,
sprzężonym procesie, stopniowo korygując oba wyniki w~miarę napływu
nowych danych.
W~zależności od tego, czy w~pamięci utrzymywana jest pełna historia pozycji
czy tylko bieżąca, rozróżnia się SLAM online i~pełny SLAM.

\subsection{Aktywny SLAM}

Aktywny SLAM rozszerza klasyczne podejście o~planowanie ruchu: robot sam
decyduje, dokąd się udać, by jak najefektywniej zmniejszyć niepewność
swojej mapy i~pozycji.
Typowym podejściem jest eksploracja sterowana zyskiem informacyjnym.
Niniejsza praca dotyczy pasywnego SLAM na gotowych nagraniach,
dlatego aktywny SLAM nie jest tu rozwijany.

\section{Sensory}

Jakość i~rodzaj dostępnych danych sensorycznych bezpośrednio determinuje,
które algorytmy SLAM mogą być zastosowane i~jakiej dokładności można się spodziewać.
W~eksperymentach użyto czterech typów sensorów.

\subsection{LiDAR}

LiDAR mierzy odległości do otaczających obiektów, emitując impulsy laserowe
i~rejestrując czas ich powrotu.
Wynikiem jest trójwymiarowa chmura punktów opisująca geometrię otoczenia
z~wysoką dokładnością i~niezależnie od warunków oświetleniowych.
Sensor nie rejestruje jednak informacji barwnej ani teksturowej,
a~przy szybkim ruchu robota chmura punktów ulega zniekształceniu geometrycznemu,
ponieważ kolejne punkty rejestrowane są w~różnych chwilach czasu.

\subsection{Kamera stereo}

Kamera stereo składa się z~dwóch obiektywów rozstawionych w~poziomie.
Porównując obrazy z~obu kamer, system wyznacza dysparycję każdego piksela,
którą przelicza na głębokość.
Im dalej od kamery znajduje się punkt, tym mniejsza jest dysparycja
i~tym większa niepewność pomiaru głębokości.
Kamery dostarczają bogatej informacji teksturowej, co jest kluczowe
dla wizyjnych algorytmów SLAM, lecz ich skuteczność spada w~słabym oświetleniu
lub środowiskach pozbawionych charakterystycznych wzorów.

\subsection{IMU}

Centrala inercyjna mierzy przyspieszenie liniowe i~prędkość kątową
z~bardzo wysoką częstotliwością, rzędu kilkuset herców.
Pozwala śledzić ruch robota między kolejnymi odczytami wolniejszych sensorów
i~korygować zniekształcenia chmur punktów LiDAR.
Wadą IMU jest dryf wynikający z~kumulacji błędów całkowania oraz biasy
pomiarowe, które muszą być estymowane i~kompensowane.

\subsection{Odometria kołowa}

Enkodery zamontowane na kołach zliczają obroty i~na tej podstawie
wyznaczany jest przyrost pozycji robota.
Metoda jest prosta i~tania, lecz podatna na poślizg kół i~błędy kalibracji
wymiarów mechanicznych.
W~systemach SLAM odometria kołowa służy zazwyczaj jako dodatkowe, niskokosztowe
źródło ruchu wspomagające dokładniejsze sensory.

\section{Reprezentacje map}

Wybór reprezentacji mapy wpływa na wymagania pamięciowe, szybkość aktualizacji
i~przydatność mapy do różnych zadań, takich jak planowanie trasy czy inspekcja.

\subsection{Siatka zajętości}

Siatka zajętości dzieli przestrzeń na jednakowe komórki, z~których każda
przechowuje szacowane prawdopodobieństwo zajętości przez przeszkodę.
Każdy odczyt sensoryczny aktualizuje odpowiednie komórki, zwiększając pewność
co do obecności lub braku przeszkód.
Reprezentacja ta jest bezpośrednio użyteczna do planowania ścieżek,
lecz ogranicza się do przestrzeni 2D i~nie koduje struktury ponad
i~pod płaszczyzną skanowania.

\subsection{Chmura punktów 3D}

Chmura punktów 3D to zbiór współrzędnych przestrzennych zarejestrowanych
przez sensor.
Mapa globalna powstaje przez łączenie kolejnych chmur z~uwzględnieniem
estymowanej pozycji robota.
Reprezentacja zachowuje pełną trójwymiarową geometrię środowiska i~nadaje się
do inspekcji oraz oceny jakości odwzorowania, lecz zajmuje dużo pamięci
i~nie zawiera explicite informacji o~wolnej przestrzeni.

\section{Rejestracja chmur punktów}

Rejestracja to proces wyznaczania przekształcenia przestrzennego, które jak
najdokładniej dopasowuje dwie zachodzące na siebie chmury punktów.
Jest podstawą skanowego dopasowania stosowanego w~odometrii LiDAR i~zamykaniu
pętli.

\subsection{Iterative Closest Point (ICP)}

ICP iteracyjnie wyznacza korespondencje między punktami obu chmur,
wyszukując dla każdego punktu chmury źródłowej jego najbliższego sąsiada
w~chmurze docelowej, a~następnie wyznacza transformację minimalizującą
sumaryczną odległość między korespondującymi parami.
Proces powtarzany jest do zbieżności.
Algorytm jest podatny na zły punkt startowy i~wymaga dostatecznego pokrycia
między chmurami.

\subsection{Point-to-Plane ICP}

Odmiana ICP, w~której minimalizowana jest odległość punktu źródłowego
od płaszczyzny stycznej do powierzchni docelowej w~punkcie korespondencji.
Dzięki uwzględnieniu lokalnej geometrii powierzchni metoda zbiega szybciej
i~jest dokładniejsza na gładkich obiektach.

\subsection{Generalized ICP (GICP)}

GICP traktuje każdy punkt jako małe, lokalne rozkłady opisujące fragment
powierzchni i~dopasowuje te rozkłady do siebie.
Łączy właściwości ICP punkt-do-punktu i~punkt-do-płaszczyzny w~jednym
spójnym sformułowaniu probabilistycznym.
Jest domyślnym wariantem w~kilku algorytmach ocenianych w~niniejszej pracy.

\section{Filtracyjna estymacja stanu}

Metody filtracyjne reprezentują stan robota jako rozkład prawdopodobieństwa
i~aktualizują go cyklicznie: predykcja propaguje stan przez model ruchu,
korekcja uwzględnia nowe pomiary.

\subsection{Rozszerzony filtr Kalmana (EKF)}

EKF zakłada, że rozkład stanu jest gaussowski i~że modele ruchu oraz obserwacji
mogą być w~każdej chwili zlinearyzowane.
W~kroku predykcji propaguje średnią i~macierz kowariancji przez zlinearyzowany
model ruchu, a~w~kroku korekcji łączy przewidywanie z~rzeczywistym pomiarem
za pomocą wzmocnienia Kalmana.
Jest podstawą wielu systemów fuzji sensorycznej, w~tym narzędzia robot\_localization
użytego w~eksperymentach.

\subsection{Filtr cząsteczkowy Rao-Blackwellized}

Filtr cząsteczkowy reprezentuje rozkład stanu jako dyskretny zbiór
próbek (\textit{cząsteczek}), z~których każda niesie hipotezę trajektorii
i~odpowiadającą jej mapę.
Po nadejściu nowego pomiaru wagi cząsteczek są aktualizowane, a~cząsteczki
o~niskich wagach zastępowane są nowymi, bliższymi bardziej prawdopodobnym
hipotezom.
Na tym mechanizmie opiera się algorytm GMapping stosowany w~pracy.

\section{Grafowy SLAM}

Grafowe podejście do SLAM dzieli problem na dwie fazy: budowanie grafu
(front-end, oparty na odometrii i~dopasowaniu skanów) i~optymalizację grafu
(back-end, zapewniającą globalną spójność mapy).

\subsection{Grafy pozycji}

Każda estymowana pozycja robota to węzeł grafu, a~każde wiązanie sensoryczne ---
krawędź z~przypisaną miarą pewności.
Optymalizacja minimalizuje sumaryczny błąd wszystkich wiązań, rozkładając
niespójności równomiernie po całej trajektorii.

\subsection{Grafy czynnikowe}

Grafy czynnikowe uogólniają grafy pozycji: węzłami mogą być dowolne zmienne
(pozy, punkty mapy, biasy IMU), a~krawędzie reprezentują probabilistyczne
zależności między nimi.
Formalizm ten jest stosowany m.in.\ przez algorytmy LIO-SAM i~MOLA.

\subsection{Detekcja i~zamykanie pętli}

Gdy robot powraca w~pobliże wcześniej odwiedzonego miejsca, system może
rozpoznać to miejsce i~dodać do grafu długodystansowe wiązanie korekcyjne.
Bez zamykania pętli błąd pozycji narasta nieograniczenie, a~mapa staje się
niespójna na długich trajektoriach.

\section{Preintegracja IMU i~deskewing}

Szybki ruch robota powoduje, że kolejne punkty jednej chmury LiDAR rejestrowane
są w~różnych chwilach czasu, co zniekształca geometrię chmury.
Korekcja tego efektu --- deskewing --- polega na przetransformowaniu każdego
punktu do wspólnej chwili z~użyciem danych IMU.

Preintegracja IMU pozwala skumulować dziesiątki lub setki pomiarów IMU
w~jedno kompaktowe wiązanie między klatkami kluczowymi, które może być
bezpośrednio włączone do grafu czynnikowego bez ponownego całkowania
przy każdej iteracji optymalizatora.

\section{Podstawy wizyjnego SLAM}

Wizyjny SLAM opiera się na obrazach kamery.
W~niniejszej pracy oceniane są metody oparte na rzadkich cechach wizyjnych.

\subsection{Stereoskopia i~triangulacja}

Kamera stereo wyznacza głębokość każdego punktu na podstawie różnicy
położenia jego obrazu w~lewym i~prawym obiektywie.
Im mniejsza ta różnica (mała dysparycja), tym większa odległość punktu
i~tym mniej pewny pomiar głębokości.

\subsection{Detekcja i~opis cech}

Cechy wizyjne to charakterystyczne fragmenty obrazu --- narożniki, plamy ---
które dają się niezawodnie wykryć i~opisać numerycznie.
Deskryptor ORB jest odporny na obroty obrazu i~szybki obliczeniowo,
co czyni go popularnym wyborem w~systemach działających w~czasie rzeczywistym.
Dopasowanie cech między klatkami odbywa się przez porównanie deskryptorów.

\subsection{Odometria wizyjna}

Dopasowane cechy między kolejnymi klatkami pozwalają wyznaczyć względną
pozycję kamery.
Lokalna optymalizacja (bundle adjustment) minimalizuje błąd reprojekcji
dla zestawu klatek i~punktów 3D, poprawiając spójność krótkoterminową.

\subsection{Rozpoznawanie miejsc metodą Bag-of-Words}

Słownik wizualny grupuje deskryptory w~klastry --- \textit{słowa wizualne}.
Każdy obraz jest wówczas reprezentowany jako histogram częstości tych słów.
Podobieństwo histogramów służy do wyszukiwania wcześniej odwiedzonych miejsc,
a~geometryczna weryfikacja przez RANSAC potwierdza poprawność kandydatów.

\section{Metryki ewaluacji}

\subsection{Wyrównanie Umeyamy}

Przed porównaniem trajektoria estymowana wyrównywana jest do referencji przez
optymalizację transformacji podobieństwa minimalizującej błąd RMSE między
odpowiadającymi sobie pozycjami.
Dla systemów metrycznych skala transformacji jest ustalona.

\subsection{Bezwzględny błąd trajektorii (ATE)}

ATE to średni kwadratowy błąd translacyjny między estymowaną a~referencyjną
pozycją robota po wyrównaniu.
Mierzy globalną spójność trajektorii i~długoterminowy dryf algorytmu.

\subsection{Względny błąd pozycji (RPE)}

RPE wyznaczany jest dla krótkich odcinków trajektorii o~stałej długości
lub czasie.
Mierzy lokalną jakość odometrii --- czyli to, jak dokładnie algorytm
odwzorowuje ruch robota w~krótkim oknie czasowym, niezależnie od
globalnego dryfu.


% ═════════════════════════════════════════════════════════════════════════════
% WERSJA C — precyzyjna, minimalna matematyka
% ═════════════════════════════════════════════════════════════════════════════

\chapter{Podstawy teoretyczne}

Rozdział zawiera terminologię i~koncepcje algorytmiczne niezbędne do
interpretacji wyników przedstawionych w~rozdziale~4.
Omówiono podproblemy nawigacji mobilnej i~ich wzajemne zależności, sensory
stosowane w~eksperymentach, reprezentacje map, metody rejestracji chmur punktów,
filtracyjne i~grafowe podejścia do estymacji stanu, preintegrację IMU,
podstawy wizyjnego SLAM oraz metryki ewaluacji.

\section{Nawigacja w~robotyce mobilnej}

Nawigacja autonomiczna polega na celowym przemieszczaniu się robota
i~tradycyjnie dekomponowana jest na trzy zadania: zliczanie drogi, mapowanie
i~lokalizację.
SLAM łączy mapowanie i~lokalizację w~jeden sprzężony problem, eliminując
konieczność posiadania gotowej mapy lub zewnętrznego systemu pozycjonowania.

\subsection{Zliczanie drogi}

Zliczanie drogi to estymacja pozy robota przez całkowanie kolejnych przyrostów
ruchu z~odometrii lub sterowań.
Nie wymaga znajomości mapy, lecz błędy kolejnych przyrostów kumulują się
bez mechanizmu korekcji, prowadząc do narastającego dryfu.
W~systemach SLAM pełni rolę modelu predykcji między kolejnymi pomiarami.

\subsection{Mapowanie}

Mapowanie to wyznaczanie przestrzennego modelu środowiska ze strumienia pomiarów
sensorycznych przy założeniu znanych poz robota.
Błąd pozycji przenosi się bezpośrednio na spójność mapy: niepoprawnie
zlokalizowane pomiary deformują globalną strukturę odwzorowania.

\subsection{Lokalizacja}

Lokalizacja to estymacja pozy robota względem danej mapy na podstawie
bieżących pomiarów.
Wyróżnia się lokalizację globalną (bez informacji apriorycznej) i~śledzenie
pozy (z~ciągłością ruchu jako aprioryczną wiedzą).

\subsection{SLAM}

SLAM szacuje jednocześnie trajektorię robota i~mapę środowiska ze strumienia
danych sensorycznych i~sterowań.
Wzajemna zależność obu zadań czyni problem trudnym: mapa poprawia lokalizację,
a~lokalizacja poprawia mapę.
Podejście online marginalizuje przeszłe pozy; podejście full SLAM zachowuje
pełną trajektorię, umożliwiając globalną optymalizację.

\subsection{Aktywny SLAM}

Aktywny SLAM rozszerza pasywne podejście o~planowanie ruchu sterowane
redukcją niepewności mapy i~trajektorii.
Niniejsza praca dotyczy pasywnego SLAM na nagranych danych; aktywny SLAM
pozostaje poza zakresem rozważań.

\section{Sensory}

Typ i~jakość dostępnych sensorów determinuje, jakie rodziny algorytmów SLAM
mogą być stosowane i~jakich charakterystyk dokładności można oczekiwać.
W~eksperymentach zastosowano cztery uzupełniające się sensory.

\subsection{LiDAR}

LiDAR mierzy odległości metodą czasu przelotu impulsu laserowego,
dostarczając trójwymiarowej chmury punktów z~zasięgiem do kilkudziesięciu metrów
i~rozdzielczością kątową rzędu dziesiątych części stopnia.
Sensor działa niezależnie od oświetlenia, lecz nie rejestruje informacji
barwnej; przy szybkim ruchu robota punkty rejestrowane podczas jednego obrotu
głowicy tworzą geometrycznie zniekształconą chmurę (\textit{motion distortion}).

\subsection{Kamera stereo}

Kamera stereo wyznacza głębokość na podstawie dysparycji między obrazami
lewego i~prawego obiektywu; dokładność maleje kwadratowo z~odległością.
Dostarcza bogatej informacji teksturowej przydatnej w~dopasowaniu cech,
lecz jest wrażliwa na słabe lub niejednolite oświetlenie i~środowiska
o~monotonnej teksturze.

\subsection{IMU}

Centrala inercyjna mierzy przyspieszenie liniowe i~prędkość kątową
z~częstotliwością rzędu $200$~Hz.
Pomiary obarczone są biasem i~szumem gaussowskim, których kumulacja prowadzi
do dryfu; ciasne sprzężenie IMU z~LiDAR-em lub kamerą pozwala jednak
efektywnie korygować zniekształcenia i~zwiększa odporność na dynamikę ruchu.

\subsection{Odometria kołowa}

Enkodery kół dostarczają przyrostów pozy przez zliczanie obrotów kół
w~modelu różnicowym.
Poślizg i~błędy kalibracji (promień koła, rozstaw kół) wprowadzają błąd
systematyczny; w~systemach SLAM odometria kołowa stosowana jest jako
niskokosztowe uzupełnienie dokładniejszych sensorów.

\section{Reprezentacje map}

\subsection{Siatka zajętości}

Siatka zajętości to dyskretna dwuwymiarowa siatka komórek przechowujących
prawdopodobieństwo zajętości w~postaci log-odds.
Każdy odczyt sensoryczny aktualizuje odpowiednie komórki za pomocą prostego
modelu obserwacji.
Reprezentacja nadaje się bezpośrednio do planowania ścieżek,
lecz ogranicza się do 2D i~ma stałą rozdzielczość.

\subsection{Chmura punktów 3D}

Chmura punktów to nieuporządkowany zbiór trójwymiarowych współrzędnych.
Mapa globalna powstaje przez akumulację kolejnych skanów transformowanych
do układu globalnego przy użyciu estymowanych poz.
Reprezentacja zachowuje pełną geometrię 3D, lecz nie koduje wolnej przestrzeni
i~jest kosztowna pamięciowo dla długich trajektorii.

\section{Rejestracja chmur punktów}

Rejestracja wyznacza transformację sztywną z~$SE(3)$ minimalizującą
geometryczną rozbieżność między dwoma zachodzącymi na siebie chmurami punktów.

\subsection{Iterative Closest Point (ICP)}

ICP iteracyjnie przypisuje korespondencje jako pary najbliższych punktów
i~wyznacza transformację minimalizującą sumę kwadratów odległości
między korespondującymi parami.
Algorytm jest wrażliwy na inicjalizację i~lokalnie zbieżny.

\subsection{Point-to-Plane ICP}

Zamiast odległości punkt-do-punktu minimalizowana jest składowa normalna
wektora błędu względem płaszczyzny stycznej do chmury docelowej.
Metryka jest bardziej informatywna na gładkich powierzchniach
i~przyspiesza zbieżność.

\subsection{Generalized ICP (GICP)}

GICP modeluje lokalne sąsiedztwo każdego punktu jako rozkład gaussowski
nad geometrią powierzchni i~minimalizuje odchylenie Mahalanobisa
między dopasowanymi rozkładami.
Unifikuje podejścia punkt-do-punktu i~punkt-do-płaszczyzny;
jest domyślnym wariantem w~kilku algorytmach ocenianych w~pracy.

\section{Filtracyjna estymacja stanu}

\subsection{Rozszerzony filtr Kalmana (EKF)}

EKF reprezentuje stan robota jako gaussowski rozkład (średnia, macierz kowariancji)
i~propaguje go przez zlinearyzowane modele ruchu i~obserwacji.
Krok predykcji przesuwa średnią i~rozszerza kowariancję zgodnie z~modelem ruchu;
krok korekcji kurczy kowariancję na podstawie nowego pomiaru proporcjonalnie
do jego wiarygodności.
Jest stosowany w~fuzji sensorycznej odometrii (robot\_localization).

\subsection{Filtr cząsteczkowy Rao-Blackwellized}

Rozkład pozy reprezentowany jest przez skończony zbiór cząsteczek,
z~których każda niesie hipotetyczną trajektorię i~warunkową mapę.
Po nadejściu pomiaru wagi cząsteczek są aktualizowane według modelu obserwacji,
a~następnie resampling zastępuje cząsteczki mało prawdopodobne.
Na tym mechanizmie opiera się GMapping.

\section{Grafowy SLAM}

\subsection{Grafy pozycji}

Trajektoria robota i~mapa reprezentowane są jako graf: węzły to estymowane
pozy, krawędzie --- wiązania z~odometrii lub dopasowania skanów z~przypisanymi
macierzami kowariancji.
Back-end minimalizuje sumę ważonych błędów wiązań przez nieliniową optymalizację
(g2o, GTSAM, Ceres), równomiernie rozkładając rezydua po całej trajektorii.

\subsection{Grafy czynnikowe}

Grafy czynnikowe uogólniają grafy pozycji: węzłami mogą być dowolne zmienne
(pozy, biasy IMU, skala), krawędziami --- dowolne czynniki probabilistyczne.
Formalizm stosowany m.in.\ przez LIO-SAM i~MOLA.

\subsection{Detekcja i~zamykanie pętli}

Zamykanie pętli wykrywa powrót robota w~pobliże wcześniej odwiedzonego miejsca
i~dodaje do grafu długodystansowe wiązanie korekcyjne.
Bez niego błąd trajektorii narasta wzdłuż całej trasy; z~nim globalny
back-end jest w~stanie zniwelować dryf nagromadzony od ostatniego zamknięcia.

\section{Preintegracja IMU i~deskewing}

Deskewing koryguje zniekształcenie geometryczne chmury punktów LiDAR
wynikające z~ruchu robota podczas obrotu głowicy: każdy punkt transformowany
jest do wspólnej chwili referencyjnej z~użyciem interpolowanej pozy z~IMU.

Preintegracja IMU łączy wszystkie pomiary między klatkami kluczowymi
w~kompaktowe wiązanie preintegrowane, reprezentowane na grupie Liego $SO(3)$,
które może być bezpośrednio włączone do grafu czynnikowego bez ponownego
całkowania przy każdej iteracji optymalizatora.

\section{Podstawy wizyjnego SLAM}

\subsection{Stereoskopia i~triangulacja}

Głębokość wyznaczana jest ze stosunku iloczynu ogniskowej i~bazy do dysparycji.
Niepewność głębokości rośnie kwadratowo z~odległością, ograniczając zakres,
w~którym pomiary stereo są wiarygodne.

\subsection{Detekcja i~opis cech}

Cechy ORB łączą detektor FAST z~binarnym deskryptorem BRIEF,
zapewniając niezmienniczość na obroty i~szybkie obliczenia.
Dopasowanie par deskryptorów filtrowane jest testem stosunku odległości
eliminującym niejednoznaczne korespondencje.

\subsection{Odometria wizyjna}

Wzajemna poza kamer wyznaczana jest z~dopasowanych korespondencji
przez dekompozycję macierzy esencjalnej lub minimalizację błędu reprojekcji
w~lokalnym bundle adjustment.
Lokalna optymalizacja poprawia krótkoterminową spójność, lecz nie zapobiega
długoterminowemu dryfowi bez zamykania pętli.

\subsection{Rozpoznawanie miejsc metodą Bag-of-Words}

Słownik wizualny dzieli przestrzeń deskryptorów na klastry --- słowa wizualne.
Każdy obraz reprezentowany jest histogramem słów (z~ważeniem TF-IDF),
a~zamknięcie pętli wyzwalane jest przez wysokie podobieństwo kosinusowe
między bieżącym a~archiwalnymi histogramami, po geometrycznej weryfikacji RANSAC.

\section{Metryki ewaluacji}

\subsection{Wyrównanie Umeyamy}

Przed obliczeniem błędów trajektoria estymowana wyrównywana jest do referencji
przez optymalizację transformacji podobieństwa minimalizującej RMSE między
odpowiadającymi pozami (rozwiązanie zamknięte SVD).
Dla systemów metrycznych skala jest ustalona na~1.

\subsection{Bezwzględny błąd trajektorii (ATE)}

ATE to RMSE błędów translacyjnych między estymowanymi a~referencyjnymi
pozycjami robota po wyrównaniu.
Mierzy globalną spójność i~długoterminowy dryf trajektorii.

\subsection{Względny błąd pozycji (RPE)}

RPE wyznaczany jest dla par poz odległych o~stały krok jako RMSE różnic
między względnymi transformacjami estymowanymi a~referencyjnymi.
Mierzy lokalną jakość odometryczną niezależnie od globalnego dryfu.
